\documentclass[letterpaper,11pt]{moderncv}
\moderncvstyle{banking}
\usepackage{xpatch}
\moderncvcolor{black}
\AtBeginDocument{\definecolor{color2}{rgb}{0,0,0}}
\usepackage[letterpaper, margin=0.6in]{geometry}
\usepackage{xcolor}
\usepackage{graphicx}
\usepackage{fontspec}
\setmainfont{Times New Roman}
\pagenumbering{arabic}
\usepackage[unicode]{hyperref}
\tolerance=1
\emergencystretch=\maxdimen
\hyphenpenalty=10000
\hbadness=10000
\usepackage{enumitem}
\usepackage{needspace}
\usepackage{lastpage}
\rfoot{\addressfont\itshape\textcolor{gray}{\thepage/\pageref{LastPage}}}

\usepackage[english]{babel}
\usepackage[autostyle, english = american]{csquotes}
\MakeOuterQuote{"}

%==================== PERSONAL INFO ====================%
\name{\textcolor{black}{Saeyoung}}{Kim}
\address{Cambridge, MA, USA}
\phone[mobile]{+1 (351) 322 5510}
\email{saeyoung\_kim@uml.edu}
\social[linkedin]{saeyoung-macx-kim}
\social[github]{PrimaryAnomaly}

\begin{document}

\makecvtitle
\vspace{-3em}

%==================== SUMMARY ====================%
\section{Professional Summary}
Technical project manager and systems engineer who has led multi-year, cross-border R\&D programs in pharmaceutical manufacturing, space hardware, and AI-driven process software. Currently managing the US side of a multi-million-dollar KEIT-funded digital twin program. Previously ran a KRW 300M (\textasciitilde\$230K) government-funded CubeSat payload project through 8 years of development. I build the software I manage: full-stack web apps, real-time data pipelines, simulation models, and multi-agent AI systems. Bilingual in English and Korean, comfortable bridging technical teams across time zones and institutions.

%==================== TECHNICAL SKILLS ====================%
\section{Technical Skills}
\begin{minipage}[t]{0.48\textwidth}
\cvitem{Project Mgmt}{WBS Development, Gantt Charts, Critical Path Analysis, Risk Tracking, Milestone Planning, Cross-Team Coordination, Government Project Reporting}
\cvitem{Software Dev}{Python, C/C++, MATLAB, FastAPI, Next.js, Docker, Git, Linux}
\end{minipage}%
\hfill
\begin{minipage}[t]{0.48\textwidth}
\cvitem{PM Tools}{Excel WBS/Gantt, PowerPoint Status Decks, Git/GitHub}
\cvitem{Data/AI}{PCA, PLS, ODE Modeling, TimescaleDB, Telegraf, Real-Time Data Pipelines}
\cvitem{Hardware}{Circuit Design, PCB Layout (KiCad), Sensor Systems, Embedded Firmware, CAD (Fusion360, SolidWorks)}
\end{minipage}

%==================== PROJECT MANAGEMENT ====================%
\section{Project Management Experience}

\cventry{2025.07--Present}{\textbf{Project Manager and Technical Lead} | Lyophilization Digital Twin}{University of Massachusetts Lowell (KEIT-funded)}{Lowell, MA, USA}{}{
\begin{itemize}[leftmargin=1cm, itemsep=0pt, parsep=0pt]
    \item Manage US side of a 5-year KEIT-funded digital twin for pharmaceutical freeze-drying (mRNA-LNP), coordinating a 6-person team with Korean partners
    \item Maintain a 15-work-package WBS with multi-year Gantt chart across two phases (3-year R\&D + 2-year technology transfer)
    \item Run weekly status meetings with executive summaries, milestone tracking, and risk/blocker identification
    \item Track procurement of specialized equipment (\$194K+ TDLAS, 16-week lead time) and coordinate vendor timelines with project milestones
    \item Assign task ownership with monthly milestones across parallel workstreams: experimentation, modeling, frontend, backend
    \item Pivoted from gPROMS to Python when the commercial tool could not support the required data pipeline, keeping the project on schedule
    \item Built the project's core software: 9-state ODE model, multi-agent AI system (DAEMON), FastAPI + Next.js web app, Telegraf/TimescaleDB pipeline
\end{itemize}}

\needspace{12\baselineskip}
\cventry{2017.03--2024.02}{\textbf{Project Manager and Technical Lead} | BioCubeSat Biological Payload (CyanOx)}{Korea University (NRF + GPF funded)}{Seoul, South Korea}{}{
\begin{itemize}[leftmargin=1cm, itemsep=0pt, parsep=0pt]
    \item Led an 8-year, KRW 300M (\textasciitilde\$230K) program across two government grants (5-year GPF + 3-year NRF) to build a 3U CubeSat biological payload
    \item Created WBS across 6 subsystems (Mechanical, Fluidics, Sensors, Control, MCU, Biology) with 12 versioned action plans and Gantt scheduling
    \item Wrote monthly work plans with per-person task assignments and collected monthly progress reports from each of 4 team members
    \item Managed the full research budget (personnel, materials, equipment, travel) with NRF/GPF compliance
    \item Delivered all milestones to the funding agency: kickoff report, 2 annual reports, 1 final report (5 to 8 revision cycles each)
    \item Ran weekly team meetings covering scope review, Gantt status, development methods, and monthly plan alignment
    \item Coordinated the design-build-test cycle across 6 subsystems through two full hardware integration rounds
    \item Published 3 first-author journal papers and delivered 4 conference presentations as project deliverables
\end{itemize}}

\needspace{8\baselineskip}
\cventry{2015--2023}{\textbf{Co-Founder and Lead Organizer} | Effective Altruism South Korea and AMF Korea}{Nonprofit Community Management}{Seoul, South Korea}{}{
\begin{itemize}[leftmargin=1cm, itemsep=0pt, parsep=0pt]
    \item Co-founded Effective Altruism South Korea in 2015, built a bilingual community with regular meetups and tracked metrics
    \item Secured over \$100K in grants from the Centre for Effective Altruism for community operations and programs
    \item Built organizational infrastructure from scratch: OKRs, member pipeline, finances, project tracking, standard processes
    \item Ran concurrent projects: book club, Korean translation of EA materials, outreach strategy, member surveys
    \item Helped found AMF Korea as a registered nonprofit in 2023, recruiting promoters and preparing bilingual legal documents
\end{itemize}}

%==================== EXPERIENCE ====================%
\needspace{6\baselineskip}
\section{Professional Experience}

\cventry{2025.07--Present}{\textbf{Postdoctoral Associate} | Manufacturing Process Monitoring and Test Automation}{University of Massachusetts Lowell}{Lowell, MA, USA}{}{
\begin{itemize}[leftmargin=1cm, itemsep=0pt, parsep=0pt]
    \item Built Python diagnostic software suite replacing commercial platforms for real-time manufacturing process monitoring
    \item Developed anomaly detection and quality prediction pipelines across 100+ pharmaceutical production batches
    \item Designed data-driven test strategies with automated error flagging for multi-day batch manufacturing cycles
    \item Translated cross-functional stakeholder requirements into standardized test procedures and documentation
\end{itemize}}

\needspace{6\baselineskip}
\cventry{2024.02--2025.05}{\textbf{Systems Integration Engineer} | Test Equipment Design and Satellite Integration}{Hanwha Systems}{Seoul, South Korea}{}{
\begin{itemize}[leftmargin=1cm, itemsep=0pt, parsep=0pt]
    \item Designed and deployed EGSE architecture for satellite subsystem factory acceptance testing
    \item Developed and executed system-level test plans for military SAR satellite qualification across EM, QM, and FM builds
    \item Delivered turnkey test setups with data acquisition hardware, sensor interfaces, and electrical connections for production
    \item Performed root cause analysis of electromechanical failures during integration, driving corrective actions to resolution
    \item Provided design-for-test recommendations based on failure trends and fault isolation during assembly campaigns
    \item Managed purchasing, integration, and commissioning of functional test systems for factory and field deployment
    \item Led test equipment reliability improvements by identifying recurring failure modes and implementing fixes
    \item Wrote and maintained controlled test procedures, fixture documentation, and troubleshooting guides
    \item Trained engineers and technicians on test equipment operation, debug processes, and troubleshooting methods
    \item Coordinated across electrical, mechanical, software, and program teams to balance test coverage with schedules
\end{itemize}}

%==================== EDUCATION ====================%
\needspace{6\baselineskip}
\section{Education}
\cventry{2017.03--2024.02}{\textbf{PhD in Electrical Engineering} | Embedded Sensor Systems and Hardware Test Development}{Korea University}{Seoul, South Korea}{}{
\begin{itemize}[leftmargin=1cm, itemsep=0pt, parsep=0pt]
    \item Secured and managed KRW 300M (\textasciitilde\$230K) in research grants for CubeSat biological payload development
    \item Built custom test fixtures and data acquisition setups using oscilloscopes, DMMs, function generators, and power supplies
    \item Designed multi-modal sensor suite through full cycle: schematics, PCB layout, fabrication, assembly, and qualification
    \item Developed MCU firmware in C/C++ for real-time multi-channel data acquisition over UART, SPI, and I2C
    \item Performed iterative environmental and reliability testing to identify failure modes and optimize fabrication parameters
    \item Diagnosed hardware/firmware interaction failures and implemented corrective actions through multiple revision cycles
    \item Published 3 first-author and 10+ co-authored journal papers on sensors, test systems, and space payload hardware
\end{itemize}}

\cventry{2013.03--2017.02}{\textbf{BSc in Mechanical Engineering}}{Konkuk University}{Seoul, South Korea}{}{
\begin{itemize}[leftmargin=1cm, itemsep=0pt, parsep=0pt]
    \item Conducted research in electrochemical sensor microfabrication, cleanroom processing, and hardware prototyping
    \item Gained hands-on experience in CAD modeling, mechanical assembly, PCB prototyping, and precision measurement
\end{itemize}}

%==================== PUBLICATIONS ====================%
\needspace{6\baselineskip}
\section{Publications}
\subsection{Journal Publications}
\begin{sloppypar}
\begin{itemize}[leftmargin=1cm, itemsep=0pt, parsep=0pt]
    \item \textbf{\underline{Saeyoung Kim}}, Jaewon Lee, Kimia Babaei, Seonkgyu Yoon, "Digital Twins in Lyophilization: Advances in Pharmaceutical Process Development and Optimization," \textit{Biotechnology Advances}, 2026. (Manuscript in Preparation, Target Submission: Q1 2026)
    \item Jinhwee Kil, \textbf{\underline{Saeyoung Kim}}, James Jungho Pak, "Investigation of thermal management strategies for biological CubeSat payloads," \textit{Acta Astronautica}, 2024, Art. no. 228.
    \item \textbf{\underline{Kim, S.}}; Park, S.; Pak, J.J., "Multi-Modal Multi-Array Electrochemical and Optical Sensor Suite for a Biological CubeSat Payload," \textit{MDPI Sensors}, 2024, 24, Art. no. 265.
    \item Ji-Hoon Han, Sang Hyun Park, \textbf{\underline{Saeyoung Kim}}, James Jungho Pak, "A Performance Improvement of Enzyme-Based Electrochemical Lactate Sensor Fabricated by Electroplating Novel PdCu Mediator on a Laser Induced Graphene Electrode," \textit{Bioelectrochemistry}, 2022, 148, 108259.
    \item \textbf{\underline{Saeyoung Kim}}, Ji-Hoon Han, Beelee Chua, James Jungho Pak, "A pH Sensing Pipette for Cross-Contamination Prevention in Industrial Fermentation," \textit{IEEE Trans. Industrial Electronics}, 2022, 69(7), 7461-7469.
    \item S.D. Raut, N.M. Shinde, B.G. Ghule, \textbf{\underline{S. Kim}}, J.J. Pak, Q. Xia, R.S. Mane, "Room-Temperature Solution Processed Sharp-Edged Nanoshapes of," \textit{Chemical Engineering Journal}, 2022, 433(2), 133627.
    \item S.E. Jeong, \textbf{\underline{S. Kim}}, J.H. Han, J. Jungho Pak, "Simple Laser-Induced Graphene Fiber Electrode Fabrication for High-Performance Heavy-Metal Sensing," \textit{Microchemical Journal}, 2022, 172, 106950.
    \item Sanghoon Cho, Jungmo Jung, \textbf{\underline{Saeyoung Kim}}, James Pak, "Conduction Mechanism and Synaptic Behaviour of Interfacial Switching AlOσ-based RRAM," \textit{Semiconductor Science and Technology}, 2020, 35(8), 085006.
    \item Ji-Hoon Han, \textbf{\underline{Saeyoung Kim}}, Jaesung Choi, Sora Kang, Youngmi Kim Pak, James Jungho Pak, "Development of Multi-Well-Based Electrochemical Dissolved Oxygen Sensor Array," \textit{Sensors and Actuators B: Chemical}, 2020, 306, 127465.
\end{itemize}
\end{sloppypar}

\subsection{Conference Presentations}
\begin{sloppypar}
\begin{itemize}[leftmargin=1cm, itemsep=0pt, parsep=0pt]
    \item \textit{Thermal Management System for a Fluidic Card in a Small Satellite Biological Payload}, Korean Institute of Electrical Engineers (KIEE) Autumn Conference, 2022.
    \item \textit{A Pipette with an Integrated pH Sensor for Liquid Sampling Applications}, Korean Institute of Electrical Engineers (KIEE) Summer Conference, 2019.
    \item \textit{IoT System for Monitoring pH, Dissolved Oxygen, Temperature, and Electrical Conductivity in Liquids}, Korean Institute of Electrical Engineers (KIEE) Summer Conference, 2018.
    \item \textit{Development Trends of Non-Invasive Glucose Sensor Technology}, Korean Institute of Electrical Engineers (KIEE) Summer Conference, 2017.
\end{itemize}
\end{sloppypar}

\subsection{Seminar Presentations}
\begin{itemize}[leftmargin=1cm, itemsep=0pt, parsep=0pt]
    \item \textit{Bioprocessing in Space: Challenges and Opportunities for Pharmaceutical Manufacturing Beyond Earth}, Lowell Center for Space Science and Technology (LoCSST), UMass Lowell, 2025.
\end{itemize}

\vfill
\centering{References available upon request.}

\end{document}
